\subsection{开映射与闭映射}

定义 1. 设 X、X' 是两个拓扑空间。映射 f: X → X' 称为开（相应地，闭）映射，如果 X 的每个开（相应地，闭）集在 f 下的像在 X' 中是开（相应地，闭）集。

特别地，此时 \(f(\mathbf{X})\) 是 \(X'\) 的开（相应地，闭）子集。

例。1) 设 A 是拓扑空间 X 的一个子空间。则典范单射 \(j: \mathbf{A} \to \mathbf{X}\) 是开（相应地，闭）映射，当且仅当 A 在 X 中是开（相应地，闭）集（\(\S\) 3，第 1 目）。

\begin{enumerate}
\def\labelenumi{\arabic{enumi})}
\setcounter{enumi}{1}
\item
  为使双射 \(f\)（从拓扑空间 \(\mathbf{X}\) 到拓扑空间 \(\mathbf{X}'\) 上的）是一个同胚，必须且只需 \(f\) 连续且开，或连续且闭。
\item
  设 \(f\) 是集合 \(X\) 到拓扑空间 \(X'\) 上的一个满射；若给 \(X\) 赋以在 \(f\) 之下 \(X'\) 的拓扑的逆像这一拓扑（§ 2，第 3 目，例 1），则 \(f\) 连续、开且闭。
\item
  在积空间
\end{enumerate}

\[
\mathbf {X} = \prod_ {i \in \mathbf {I}} \mathbf {X} _ {i},
\]

中，每个投影 \(\mathbf{pr}_{t}:\mathbf{X}\to\mathbf{X}_{t}\) 都是连续的开映射，但未必是闭映射（§ 4，第 2 目，命题 5）。

* 5) 在 C 的开子集 A 上的全纯函数是 A 到 C 中的一个开映射。*

\begin{enumerate}
\def\labelenumi{\arabic{enumi})}
\setcounter{enumi}{5}
\tightlist
\item
  设 X、X' 是两个拓扑空间，f 是 X 到 X' 上的一个连续但非双连续的双射。则逆双射 g: \(X' \rightarrow X\) 是 \(X'\) 到 X 上的一个既开又闭的映射，但不连续。
\end{enumerate}

命题 1. 设 X、X'、X'' 是三个拓扑空间，并设 f: X → X'、g : X' → X'' 是两个映射。则：

\begin{enumerate}
\def\labelenumi{\alph{enumi})}
\item
  若 \(f\) 与 \(g\) 是开（相应地，闭）映射，则 \(g \circ f\) 亦然。
\item
  若 \(g \circ f\) 是开（相应地，闭）映射且 \(f\) 连续并满，则 \(g\) 是开（相应地，闭）映射。
\item
  若 \(g \circ f\) 是开（相应地，闭）映射且 \(g\) 连续并单，则 \(f\) 是开（相应地，闭）映射。
\end{enumerate}

a) 由定义 1 立即得出。为证 b)，只须注意，每个开（相应地，闭）子集 \(\mathbf{A}'\)（\(\mathbf{X}'\) 的）都可写成 \(f(\mathbf{A})\)，其中 \(\mathbf{A} = \overline{f}^{1}(\mathbf{A}')\) 在 \(\mathbf{X}\) 中是开（相应地，闭）集（\(\S 2\)，第 1 目，定理 1）；从而 \(g(\mathbf{A}') = g(f(\mathbf{A}))\) 在 \(\mathbf{X}''\) 中是开（相应地，闭）集。最后，为证 c)，我们注意，\(f(\mathbf{A}) = \overline{g}^{1}(g(f(\mathbf{A}))\) 对每个子集 \(\mathbf{A}\)（\(\mathbf{X}\) 的）成立；按假设，若 \(\mathbf{A}\) 在 \(\mathbf{X}\) 中是开（相应地，闭）集，则 \(g(f(\mathbf{A}))\) 在 \(\mathbf{X}''\) 中是开（相应地，闭）集，从而 \(f(\mathbf{A})\) 在 \(\mathbf{X}'\) 中是开（相应地，闭）集（\(\S 2\)，第 1 目，定理 1）。

命题 2. 设 X、Y 是两个拓扑空间，f 是 X 到 Y 中的一个映射。对 Y 的每个子集 T，用 \(f_{T}\) 表示 \(\overline{f}^{1}(T)\) 到 T 中的、在 \(\overline{f}^{1}(T)\) 上与 f 一致的映射。

\begin{enumerate}
\def\labelenumi{\alph{enumi})}
\item
  若 \(f\) 是开（相应地，闭）映射，则 \(f_{\mathbf{T}}\) 是开（相应地，闭）映射。
\item
  设 \((\mathbf{T}(\iota))_{\iota \in \mathbf{I}}\) 是 Y 的子集的一个族，其诸内部覆盖 Y，或者它是 Y 的一个局部有限闭覆盖；若所有 \(f_{\mathbf{T}(\iota)}\) 都是开（相应地，闭）映射，则 \(f\) 是开（相应地，闭）映射。
\item
  若 A 是 \(\overline{f}^{1}(\mathbf{T})\) 的开（相应地，闭）子集，则存在 X 的开（相应地，闭）子集 B，使 \(\mathbf{A} = \mathbf{B} \cap \overline{f}^{1}(\mathbf{T})\)，从而 \(f_{\mathbf{T}}(\mathbf{A}) = f(\mathbf{B}) \cap \mathbf{T}\)；按假设，\(f(\mathbf{B})\) 是开（相应地，闭）集，于是 \(f_{\mathbf{T}}(\mathbf{A})\) 在 \(\mathbf{T}\) 中是开（相应地，闭）集。
\item
  设 B 是 X 的开（相应地，闭）子集，并用 \(B_{i}\) 表示 \(\mathbf{B} \cap \overline{f}^{1}(\mathrm{T}(i))\)；则 \(f(\mathbf{B}) \cap \mathrm{T}(i) = f_{\mathbf{T}(i)}(\mathbf{B}_{i})\)。由于按假设 \(f_{\mathbf{T}(i)}(\mathbf{B}_{i})\) 在 \(\mathrm{T}(i)\) 中是开（相应地，闭）集，故由 § 3 第 1 目命题 3，\(f(\mathbf{B})\) 在 Y 中是开（相应地，闭）集。
\end{enumerate}

推论. 设 \((\mathbf{T}(\iota))_{\iota \in \mathbf{I}}\) 是 Y 的子集的一个族，其诸内部覆盖 Y，或者它是 Y 的一个局部有限闭覆盖。若 \(f: \mathbf{X} \to \mathbf{Y}\) 连续，且每个 \(f_{\mathbf{T}(\iota)}\) 都是 \(\overline{f}^1 (\mathbf{T}(\iota))\) 到 \(\mathbf{T}(\iota)\) 上的一个同胚，则 \(f\) 是 \(\mathbf{X}\) 到 \(\mathbf{Y}\) 上的一个同胚。

因为 \(f\) 显然是双射，且凭借命题 2 是开映射。

